\documentclass[10pt,a4paper]{article}
\usepackage{amsmath,amssymb,amsthm,graphicx}
\newcommand{\Corr}{\operatorname{Corr}}
\newcommand{\DT}{\operatorname{DTDepth}}
\title{Exponential XOR hardness amplification through random simplification}
\author{William M. Hoza}
\date{doi:10.4230/LIPIcs.CCC.2024.1; CC BY 4.0}
\begin{document}
\maketitle
\noindent\textit{Allowed standard prerequisites (editorial).}\begin{itemize}
\item Decision trees, conditional expectation and product distributions. Author Definitions 2.1, 2.4, 2.5; necessary definitions and full author Claim3.3 retained.
\end{itemize}\medskip
\section*{Definitions and notation (author Section 2)}
We write $\mathbb N$ to denote the set of non-negative integers.
In the introduction, we worked with functions $f:\{0,1\}^n\to\{0,1\}$. Going forward, it will be more convenient to encode a bit $b\in\{0,1\}$ as the value $(-1)^b$. Thus, we will work with functions $f:\{\pm1\}^n\to\{\pm1\}$. We continue to use the notation $f^{\oplus t}$, but now $f^{\oplus t}$ denotes the product of $t$ copies of $f$ on independent inputs.
We use the following notation to describe decision trees.
\paragraph{Definition 2.1 (Decision trees).} For a function $f:\{\pm1\}^n\to\{\pm1\}$, we define $\DT(f)$ to be the minimum depth of a decision tree computing $f$. In the other direction, for a parameter $D\in\mathbb N$, we define $\DT[D]$ to be the class of all functions $f:\{\pm1\}^n\to\{\pm1\}$ that can be computed by depth-$D$ decision trees. (The parameter $n$ will always be clear from context.)
We denote random variables using boldface. We write $\mathbf x\sim\mu$ to indicate that the random variable $\mathbf x$ is sampled from the distribution $\mu$.
\paragraph{Definition 2.4 (Tensor product of probability distributions).} Let $\mu_1,\ldots,\mu_t$ be probability distributions over the spaces $\Omega_1,\ldots,\Omega_t$. Sample $\mathbf x^1\sim\mu_1,\ldots,\mathbf x^t\sim\mu_t$ independently. The tensor product $\mu_1\otimes\cdots\otimes\mu_t$ is the probability distribution of $(\mathbf x^1,\ldots,\mathbf x^t)$. As a special case, we define
\[\mu^{\otimes t}=\underbrace{\mu\otimes\mu\otimes\cdots\otimes\mu}_{t\text{ copies}}.\]
We use the following standard definition to reason about average-case hardness of $\{\pm1\}$-valued functions.
\paragraph{Definition 2.5 (Correlation).} Let $g,h:\{\pm1\}^n\to\mathbb R$ be functions and let $\mu$ be a distribution over $\{\pm1\}^n$. We define
\[\Corr_\mu(g,h)=\mathbb E_{\mathbf x\sim\mu}[g(\mathbf x)\cdot h(\mathbf x)].\]
More generally, if $\mathcal C$ is a class of functions $g:\{\pm1\}^n\to\mathbb R$, then we define
\[\Corr_\mu(\mathcal C,h)=\max_{g\in\mathcal C}\Corr_\mu(g,h).\]
If $\mu$ is omitted, then by default it is assumed to be the uniform distribution over $\{\pm1\}^n$.
If $g$ and $h$ are $\{\pm1\}$-valued, then a bound $|\Corr(g,h)|\le\varepsilon$ is equivalent to the statement that $g$ agrees with $h$ on at most a $(1/2+\varepsilon/2)$-fraction of inputs, because for any two $\{0,1\}$-valued random variables $\mathbf a,\mathbf b$, we have $\Pr[\mathbf a=\mathbf b]=\frac12+\frac12\mathbb E[(-1)^{\mathbf a}\cdot(-1)^{\mathbf b}]$.
\paragraph{Definition 2.9 (Log-concave).} A function $f:[0,\infty)\to(0,\infty)$ is log-concave if $\log f$ is concave, i.e., for every $x,y\in[0,\infty)$ and $\lambda\in(0,1)$, we have $f(x)^\lambda\cdot f(y)^{1-\lambda}\le f(\lambda x+(1-\lambda)y)$.
If $f$ is log-concave, then by induction on $t$, we have $\prod_{i=1}^t f(x_i)\le f(\bar x)^t$ where $\bar x=\frac1t\sum_{i=1}^t x_i$.

\paragraph{Definition 2.2 (Max-divergence).} Let $\mu$ and $\widetilde\mu$ be discrete probability distributions over some set $\Omega$. The max-divergence of $\widetilde\mu$ from $\mu$ is defined by
\[D_\infty(\widetilde\mu\Vert\mu)=\ln\left(\max_{x\in\Omega}\left(\frac{\Pr[\widetilde{\mathbf x}=x]}{\Pr[\mathbf x=x]}\right)\right),\]
where $\mathbf x\sim\mu$ and $\widetilde{\mathbf x}\sim\widetilde\mu$.
Max-divergence and total variation distance are related by the following lemma.
\paragraph{Lemma 2.3 (Low max-divergence $\Rightarrow$ low total variation distance).} Let $\mu$ and $\widetilde\mu$ be discrete probability distributions over the same set $\Omega$. Let $\varepsilon=D_\infty(\widetilde\mu\Vert\mu)$. There exists a probability distribution $\mu'$ such that $\mu$ can be written as a convex combination $\mu=(1-\varepsilon)\cdot\widetilde\mu+\varepsilon\cdot\mu'$. Moreover, the total variation distance between $\mu$ and $\widetilde\mu$ is at most $\varepsilon$.
\begin{proof}
If $\varepsilon=0$, the lemma is trivial, so assume $\varepsilon>0$. For each $x\in\Omega$, define
\[p(x)=\frac{\Pr[\mathbf x=x]-(1-\varepsilon)\Pr[\widetilde{\mathbf x}=x]}{\varepsilon},\]
where $\mathbf x\sim\mu$ and $\widetilde{\mathbf x}\sim\widetilde\mu$. Then $\sum_{x\in\Omega}p(x)=1$. Furthermore, $p(x)\ge0$, because
\[(1-\varepsilon)\Pr[\widetilde{\mathbf x}=x]\le(1-\varepsilon)\cdot e^\varepsilon\cdot\Pr[\mathbf x=x]\le\Pr[\mathbf x=x].\]
Therefore, $p(\cdot)$ is a probability mass function, and we can let $\mu'$ be the corresponding probability distribution. For the ``moreover'' part, observe that for any $S\subseteq\Omega$, we have
\[\Pr[\mathbf x\in S]=(1-\varepsilon)\cdot\Pr[\widetilde{\mathbf x}\in S]+\varepsilon\cdot\Pr[\mathbf x'\in S],\]
where $\mathbf x'\sim\mu'$. Therefore,
\[\Pr[\mathbf x\in S]\le\Pr[\widetilde{\mathbf x}\in S]+\varepsilon\cdot\Pr[\mathbf x'\in S]\le\Pr[\widetilde{\mathbf x}\in S]+\varepsilon,\]
and
\[\Pr[\mathbf x\in S]\ge(1-\varepsilon)\cdot\Pr[\widetilde{\mathbf x}\in S]\ge\Pr[\widetilde{\mathbf x}\in S]-\varepsilon.\]
\end{proof}
\paragraph{Definition 2.6 (Generalized restriction).} A generalized restriction is a function $\pi:\{\pm1\}^r\to\{\pm1\}^n$. If $f:\{\pm1\}^n\to\{\pm1\}$ is a Boolean function, then we define $g|_\pi$ to be the composition $g\circ\pi$. That is, $g|_\pi:\{\pm1\}^r\to\{\pm1\}$ is given by $g|_\pi(x)=g(\pi(x))$.
Traditional restrictions can be viewed as a special case of generalized restrictions as follows.
\paragraph{Definition 2.7 (Traditional restrictions as generalized restrictions).} A restriction is a string $\rho\in\{+1,-1,\star\}^n$. We identify $\rho$ with a generalized restriction $\pi:\{\pm1\}^r\to\{\pm1\}^n$, where $r=|\rho^{-1}(\star)|$, as follows. Given $y\in\{\pm1\}^r$, we let $\pi(y)$ be $\rho$, except that the $i$-th star is replaced with $y_i$ for every $i\in[r]$.
Next, we consider distributions over generalized restrictions, and we explain how to interpret the tensor product of such distributions.
\paragraph{Definition 2.8 (Tensor product of generalized restriction distributions).} Let $r,n\in\mathbb N$, and let $\mathcal R$ be a distribution over generalized restrictions $\pi:\{\pm1\}^r\to\{\pm1\}^n$. Let $\boldsymbol\pi_1,\ldots,\boldsymbol\pi_t$ be independent samples from $\mathcal R$, and define $\vec{\boldsymbol\pi}:\{\pm1\}^{rt}\to\{\pm1\}^{nt}$ by concatenating, i.e.,
\[\vec{\boldsymbol\pi}(y^{(1)},\ldots,y^{(t)})=(\boldsymbol\pi_1(y^{(1)}),\ldots,\boldsymbol\pi_t(y^{(t)})).\]
Then the tensor product $\mathcal R^{\otimes t}$ is the distribution of the random variable $\vec{\boldsymbol\pi}$.

\paragraph{Lemma 4.2 (Tighter XOR lemma for the random simplification method).}
Let $n,t\in\mathbb N$ and let $h:\{\pm1\}^n\to\{\pm1\}$ be a Boolean function. Let $\mathcal C$ be a class of Boolean functions $g:\{\pm1\}^{nt}\to\{\pm1\}$ that is closed under restrictions.\footnote{Every function in $\mathcal C$ has domain $\{\pm1\}^{nt}$. When we say that $\mathcal C$ is closed under restrictions, we are thinking of a restriction of $g\in\mathcal C$ as another function on $nt$ bits that ignores some of its input variables.} Let $r\in\mathbb N$, let $\mathcal R$ be a distribution over generalized restrictions $\pi:\{\pm1\}^r\to\{\pm1\}^n$, and let $\mu$ be a distribution over $\{\pm1\}^r$. Let $\varepsilon>0$, and assume the following.
\begin{enumerate}
\item (The distribution $\mu$ approximately completes $\mathcal R$ to the uniform distribution.) If we sample $\boldsymbol\pi\sim\mathcal R$ and $\mathbf y\sim\mu$ independently, and we sample $\mathbf x\in\{\pm1\}^n$ uniformly at random, then $D_\infty(\boldsymbol\pi(\mathbf y)\Vert\mathbf x)\le\varepsilon$.
\item (The class $\mathcal C$ simplifies under $\mathcal R^{\otimes t}$.) For every $g\in\mathcal C$ and every $D\in\mathbb N$, we have
\[\Pr_{\vec{\boldsymbol\pi}\sim\mathcal R^{\otimes t}}[\DT(g|_{\vec{\boldsymbol\pi}})\ge D]\le2^{t-D}.\]
\item (The function $h$ retains structure under $\mathcal R$.) For every $D\in\mathbb N$ and every $\pi\in\operatorname{Supp}(\mathcal R)$, we have
\[\Corr_\mu(h|_\pi,\DT[D])\le\varepsilon\cdot2^{D/3}.\]
\end{enumerate}
Then $\Corr(\mathcal C,h^{\oplus t})\le O(\varepsilon)^t$.

\paragraph{Lemma 3.2 (XOR lemma for decision trees under robust hardness assumptions, general version).}
Let $h_1,\ldots,h_t:\{\pm1\}^r\to\{\pm1\}$ be functions, and define $h(y^{(1)},\ldots,y^{(t)})=\prod_{i=1}^t h_i(y^{(i)})$. Let $\mu_1,\ldots,\mu_t$ be distributions over $\{\pm1\}^r$, and define $\mu=\mu_1\otimes\cdots\otimes\mu_t$. Let $\varepsilon:[0,\infty)\to(0,\infty)$ be a log-concave function, and assume that for every $i\in[t]$ and every $D\in\mathbb N$, we have
\[\Corr_{\mu_i}(h_i,\DT[D])\le\varepsilon(D).\]
Then for every $D\in\mathbb N$, we have
\[\Corr_\mu(h,\DT[Dt/2])\le O(\varepsilon(D))^t.\]

The first step of the proof of Lemma 3.2 is the following claim, which enables us to relate the success probability of a tree to the success probabilities of its subtrees.
\paragraph{Claim 3.3 (Law of total correlation).} Let $h,T,E:\{\pm1\}^r\to\{\pm1\}$. Let $\mu$ be a distribution over $\{0,1\}^r$. For each $b\in\{\pm1\}$, let $p_b=\Pr_{\mathbf y\sim\mu}[E(\mathbf y)=b]$, and let $\mu_b$ be the conditional distribution $(\mathbf y\sim\mu\mid E(\mathbf y)=b)$. Suppose that $T$ can be decomposed in the form
\[T(y)=\begin{cases}T_{+1}(y)&\text{if }E(y)=+1,\\T_{-1}(y)&\text{if }E(y)=-1\end{cases}\]
for some $T_{+1},T_{-1}:\{\pm1\}^r\to\{\pm1\}$. Then
\[\Corr_\mu(h,T)=\sum_{b\in\{\pm1\}}p_b\cdot\Corr_{\mu_b}(h,T_b).\]
\begin{proof}
\begin{align*}
\Corr_\mu(h,T)&=\mathbb E_{\mathbf y\sim\mu}[h(\mathbf y)\cdot T(\mathbf y)]\\
&=\sum_{b\in\{\pm1\}}p_b\cdot\mathbb E_{\mathbf y\sim\mu}[h(\mathbf y)\cdot T(\mathbf y)\mid E(\mathbf y)=b] &&\text{(Law of total expectation)}\\
&=\sum_{b\in\{\pm1\}}p_b\mathbb E_{\mathbf y\sim\mu_b}[h(\mathbf y)\cdot T_b(\mathbf y)].
\end{align*}
\end{proof}
Next, we consider the following notion of ``fair'' decision trees due to Shaltiel [61].
\paragraph{Definition 3.4 ($(D_1,\ldots,D_t)$-fair decision trees [61]).} Let $T:\{\pm1\}^{rt}\to\{\pm1\}$ be a decision tree and let $D_1,\ldots,D_t\in\mathbb N$. We say that $T$ is $(D_1,\ldots,D_t)$-fair if for every input $\vec y=(y^{(1)},\ldots,y^{(t)})\in(\{\pm1\}^r)^t$, for every $i\in[t]$, the computation $T(\vec y)$ makes at most $D_i$ queries to $y^{(i)}$.
The key to proving Lemma 3.2 is to generalize Definition 3.4 to the case of a set of tuples $(D_1,\ldots,D_t)$.
\paragraph{Definition 3.5 ($Q$-fair decision trees).} Let $T:\{\pm1\}^{rt}\to\{\pm1\}$ be a decision tree and let $Q\subseteq\mathbb N^t$. We say that $T$ is $Q$-fair if for every input $\vec y=(y^{(1)},\ldots,y^{(t)})\in(\{\pm1\}^r)^t$, there is some tuple $(D_1,\ldots,D_t)\in Q$ such that for every $i\in[t]$, the computation $T(\vec y)$ makes at most $D_i$ queries to $y^{(i)}$.
We emphasize that the tuple $(D_1,\ldots,D_t)$ is permitted to vary from one input $\vec y$ to another. Therefore, the fact that a tree is $Q$-fair does not necessarily imply that there is some $(D_1,\ldots,D_t)\in Q$ such that the tree is $(D_1,\ldots,D_t)$-fair. Given the concept of $Q$-fairness, it is relatively straightforward to prove the following claim by induction on the depth of $T$. The claim generalizes the analysis by Shaltiel [61], who considered the case of $(D_1,\ldots,D_t)$-fair decision trees and focused on the uniform distribution.
\paragraph{Claim 3.6 (XOR lemma for $Q$-fair decision trees).} Let $h_1,\ldots,h_t:\{\pm1\}^r\to\{\pm1\}$ be functions, and define $h(y^{(1)},\ldots,y^{(t)})=\prod_{i=1}^t h_i(y^{(i)})$. Let $\mu_1,\ldots,\mu_t$ be distributions over $\{\pm1\}^r$, and define $\mu=\mu_1\otimes\cdots\otimes\mu_t$. Let $Q\subseteq\mathbb N^t$ and let $T:\{\pm1\}^{rt}\to\{\pm1\}$ be a $Q$-fair decision tree. Then
\[\Corr_\mu(h,T)\le\sum_{(D_1,\ldots,D_t)\in Q}\prod_{i=1}^t\Corr_{\mu_i}(h_i,\DT[D_i]).\]
\begin{proof}
Assume without loss of generality that $T$ never queries the same variable twice. For the base case, if $T$ has depth 0, then $T$ is a constant function, so
\[|\Corr_\mu(h,T)|=\prod_{i=1}^t\left|\mathbb E_{\mathbf y^{(i)}\sim\mu_i}[h_i(\mathbf y^{(i)})]\right|=\prod_{i=1}^t\Corr_{\mu_i}(h_i,\DT[0]).\]
Since $T$ is $Q$-fair, $Q$ must be nonempty. The lemma follows because $\Corr_{\mu_i}(h_i,\DT[0])\le\Corr_{\mu_i}(h_i,\DT[D_i])$ for every $D_i\in\mathbb N$. For the inductive step, let $y_{j^*}^{(i^*)}$ be the variable queried by the root of the tree. Let $T_{+1}$ and $T_{-1}$ be the children of the root, corresponding to the cases $y_{j^*}^{(i^*)}=+1$ and $y_{j^*}^{(i^*)}=-1$ respectively. Define
\[Q'=\{(D_1,\ldots,D_{i^*-1},D_{i^*}-1,D_{i^*+1},\ldots,D_t):(D_1,\ldots,D_t)\in Q\text{ and }D_{i^*}\ne0\}.\]
Then $T_{+1}$ and $T_{-1}$ are both $Q'$-fair.
For each $b\in\{\pm1\}$, define
\[p_b=\Pr_{\mathbf y^{(i^*)}\sim\mu_{i^*}}[\mathbf y_{j^*}^{(i^*)}=b].\]
Let $\mu_{i^*}^b$ be the conditional distribution $(\mathbf y^{(i^*)}\sim\mu_{i^*}\mid\mathbf y_{j^*}^{(i^*)}=b)$, and for $i\ne i^*$, let $\mu_i^b=\mu_i$. Let $\mu^b=\mu_1^b\otimes\cdots\otimes\mu_t^b$. By Claim 3.3 and the induction hypothesis, we have
\begin{align*}
\Corr_\mu(h,T)&=\sum_{b\in\{\pm1\}}p_b\cdot\Corr_{\mu^b}(h,T_b)\\
&\le\sum_{b\in\{\pm1\}}p_b\cdot\sum_{(D_1,\ldots,D_t)\in Q'}\prod_{i=1}^t\Corr_{\mu_i^b}(h_i,\DT[D_i])\\
&=\sum_{(D_1,\ldots,D_t)\in Q'}\left(\sum_{b\in\{\pm1\}}p_b\cdot\Corr_{\mu_{i^*}^b}(h_{i^*},\DT[D_{i^*}])\right)\cdot\Pi_{\ne i^*}(D_1,\ldots,D_t)
\end{align*}
where $\Pi_{\ne i^*}(D_1,\ldots,D_t)=\prod_{i\in[t],i\ne i^*}\Corr_{\mu_i}(h_i,\DT[D_i])$. Now we bound the inner sum. By Claim 3.3, for any $D_{i^*}$, we have
\[\Corr_{\mu_{i^*}}(h_{i^*},\DT[D_{i^*}+1])\ge\sum_{b\in\{\pm1\}}p_b\cdot\Corr_{\mu_{i^*}^b}(h_{i^*},\DT[D_{i^*}]),\]
because we can approximate $h_{i^*}$ with respect to $\mu_{i^*}$ by first querying $y_{j^*}^{(i^*)}$ and then using optimal subtrees of depth $D_{i^*}$. For every $(D_1,\ldots,D_t)\in Q'$, we have $(D_1,\ldots,D_{i^*-1},D_{i^*}+1,D_{i^*+1},\ldots,D_t)\in Q$. Therefore,
\[\Corr_\mu(h,T)\le\sum_{(D_1,\ldots,D_t)\in Q}\prod_{i=1}^t\Corr_{\mu_i}(h_i,\DT[D_i]).\]
\end{proof}
Given Claim 3.6, our XOR lemma for decision trees under a robust hardness assumption (Lemma 3.2) readily follows, as we now show.
\begin{proof}[Proof of Lemma 3.2]
Let $T:\{\pm1\}^{rt}\to\{\pm1\}$ be a decision tree of depth at most $Dt/2$. Let $Q$ be the set of $t$-tuples $(D_1,\ldots,D_t)\in\mathbb N^t$ such that (1) $D_1+\cdots+D_t\le Dt$ and (2) $D_i$ is an integer multiple of $\lceil D/2\rceil$ for every $i$. We claim that $T$ is $Q$-fair. Indeed, let $\vec y=(y^{(1)},\ldots,y^{(t)})$ be any input, and let $D_i$ be the number of queries that $T(\vec y)$ makes to $y^{(i)}$. Let $D_i'$ be the smallest integer multiple of $\lceil D/2\rceil$ such that $D_i\le D_i'$. Then $D_i'\le D_i+(\lceil D/2\rceil-1)$, and hence $D_1'+\cdots+D_t'\le Dt/2+t\cdot(\lceil D/2\rceil-1)\le Dt$, showing that $(D_1',\ldots,D_t')\in Q$.
Therefore, by Claim 3.6,
\[\Corr_\mu(h,T)\le\sum_{(D_1,\ldots,D_t)\in Q}\prod_{i=1}^t\Corr_{\mu_i}(h_i,\DT[D_i]).\]
For any $(D_1,\ldots,D_t)\in Q$, we can define $(D_1',\ldots,D_t')$ such that $D_i'\ge D_i$ and $D_1'+\cdots+D_t'$ is exactly $Dt$ rather than being at most $Dt$. Then $\Corr_{\mu_i}(h_i,\DT[D_i])\le\Corr_{\mu_i}(h_i,\DT[D_i'])$, so
\begin{align*}
\Corr_\mu(h,T)&\le\sum_{(D_1,\ldots,D_t)\in Q}\prod_{i=1}^t\Corr_{\mu_i}(h_i,\DT[D_i'])\\
&\le\sum_{(D_1,\ldots,D_t)\in Q}\prod_{i=1}^t\varepsilon(D_i')\\
&\le\sum_{(D_1,\ldots,D_t)\in Q}\varepsilon(D)^t&&\text{(Log-concavity)}\\
&=|Q|\cdot\varepsilon(D)^t.
\end{align*}
To bound $|Q|$, observe that if $(D_1,\ldots,D_t)\in Q$, then we can write $D_i=c_i\cdot\lceil D/2\rceil$ for some nonnegative integers $c_1,\ldots,c_t$. Furthermore, $Dt\ge\sum_i c_i\cdot\lceil D/2\rceil\ge(D/2)\cdot\sum_i c_i$, so $c_1+\cdots+c_t\le2t$. Therefore, $|Q|$ is at most the number of ways that $2t$ can be partitioned into $t+1$ nonnegative integers, which is precisely $\binom{3t}{t}$. Thus,
\[\Corr_\mu(h,T)\le\binom{3t}{t}\cdot\varepsilon(D)^t\le O(\varepsilon(D))^t.\]
\end{proof}

\begin{proof}
Fix any $g\in\mathcal C$. Our job is to analyze the correlation between $g$ and $h^{\oplus t}$ under a uniform random input. By Lemma 2.3, we can sample a uniform random input by the following procedure.
\begin{enumerate}
\item Sample $\vec{\boldsymbol\pi}=(\boldsymbol\pi_1,\ldots,\boldsymbol\pi_t)\sim\mathcal R^{\otimes t}$.
\item Sample $\vec{\mathbf y}=(\mathbf y^{(1)},\ldots,\mathbf y^{(t)})\sim\mu^{\otimes t}$.
\item Sample $\vec{\mathbf e}=(\mathbf e^{(1)},\ldots,\mathbf e^{(t)})\sim(\mu')^{\otimes t}$, where $\mu'$ is the distribution over $\{\pm1\}^n$ from Lemma 2.3.
\item Sample $\mathbf I\subseteq[t]$ where $\Pr[i\in\mathbf I]=1-\varepsilon$ independently for every $i$.
\item Output the string $\vec{\mathbf x}=(\mathbf x^{(1)},\ldots,\mathbf x^{(t)})\in\{\pm1\}^{nt}$, where
\[\mathbf x^{(i)}=\begin{cases}\boldsymbol\pi_i(\mathbf y^{(i)})&\text{if }i\in\mathbf I\\\mathbf e^{(i)}&\text{if }i\notin\mathbf I.\end{cases}\]
\end{enumerate}
Let $\mathbf g:\{\pm1\}^{nt}\to\{\pm1\}$ be the function obtained from $g$ by plugging $\mathbf e^{(i)}$ into each block $i\notin\mathbf I$ and leaving the blocks in $\mathbf I$ alive. Since $\mathbf g$ ignores the variables in blocks outside $\mathbf I$, we have
\[g(\vec{\mathbf x})=\mathbf g|_{\vec{\boldsymbol\pi}}(\vec{\mathbf y}).\]
Similarly, define $\mathbf h:\{\pm1\}^{nt}\to\{\pm1\}$ by the formula
\[\mathbf h(\vec x)=\left(\prod_{i\in\mathbf I}h_i(x^{(i)})\right)\cdot\left(\prod_{i\notin\mathbf I}h_i(\mathbf e^{(i)})\right),\]
so that $h^{\oplus t}(\vec{\mathbf x})=\mathbf h|_{\vec{\boldsymbol\pi}}(\vec{\mathbf y})$. That way,
\[\Corr(g,h^{\oplus t})=\mathbb E[g(\vec{\mathbf x})\cdot h^{\oplus t}(\vec{\mathbf x})]=\mathbb E_{\mathbf I,\vec{\mathbf e},\vec{\boldsymbol\pi}}[\Corr_{\mu^{\otimes t}}(\mathbf g|_{\vec{\boldsymbol\pi}},\mathbf h|_{\vec{\boldsymbol\pi}})].\]
Let $\mathbf D=\lfloor\DT(\mathbf g|_{\vec{\boldsymbol\pi}})/|\mathbf I|\rfloor$. Then
\[\mathbb E_{\mathbf I,\vec{\mathbf e},\vec{\boldsymbol\pi}}[\Corr_{\mu^{\otimes t}}(\mathbf g|_{\vec{\boldsymbol\pi}},\mathbf h|_{\vec{\boldsymbol\pi}})]\le\mathbb E_{\mathbf I,\vec{\mathbf e},\vec{\boldsymbol\pi}}[\Corr_{\mu^{\otimes t}}(\mathbf h|_{\vec{\boldsymbol\pi}},\DT[(\mathbf D+1)\cdot|\mathbf I|])].\]
Let $\boldsymbol\pi_{\mathbf I}=(\boldsymbol\pi_i)_{i\in\mathbf I}$, and define $h_I:\{\pm1\}^{n|I|}\to\{\pm1\}$ by $h_I((x^{(i)})_{i\in I})=\prod_{i\in I}h(x^{(i)})$. Then for any fixing of $\mathbf I,\vec{\mathbf e},\vec{\boldsymbol\pi}$, we have
\[\Corr_{\mu^{\otimes t}}(\mathbf h|_{\vec{\boldsymbol\pi}},\DT[(\mathbf D+1)\cdot|\mathbf I|])=\Corr_{\mu^{\otimes|\mathbf I|}}(h_{\mathbf I}|_{\boldsymbol\pi_{\mathbf I}},\DT[(\mathbf D+1)\cdot|\mathbf I|]).\]
Now we apply Lemma 3.2. For each $i\in\mathbf I$ and each $D\in\mathbb N$, we have
\[\Corr_\mu(h|_{\boldsymbol\pi_i},\DT[D])\le\varepsilon\cdot2^{D/3}.\]
Furthermore, the function $\varepsilon(D)=\varepsilon\cdot2^{D/3}$ is log-concave. Therefore, Lemma 3.2 guarantees that
\[\Corr_{\mu^{\otimes|\mathbf I|}}(h_{\mathbf I}|_{\boldsymbol\pi_{\mathbf I}},\DT[(\mathbf D+1)\cdot|\mathbf I|])\le O(\varepsilon\cdot2^{2\cdot(\mathbf D+1)/3})^{|\mathbf I|}=O(\varepsilon\cdot2^{2\mathbf D/3})^{|\mathbf I|}.\]
Thus, overall, we get
\[\Corr(g,h^{\oplus t})\le\mathbb E_{\mathbf I,\vec{\mathbf e},\vec{\boldsymbol\pi}}\left[O\left(\varepsilon\cdot2^{2\mathbf D/3}\right)^{|\mathbf I|}\right].\]
Now consider any fixing of $\mathbf I$ and $\vec{\mathbf e}$. Since $\mathcal C$ is closed under restrictions, $\mathbf g\in\mathcal C$. Therefore, our simplification assumption tells us that for every $D\in\mathbb N$, we have
\[\Pr_{\vec{\boldsymbol\pi}}[\mathbf D=D]\le2^{t-D\cdot|\mathbf I|}.\]
Consequently,
\begin{align*}
\mathbb E_{\vec{\boldsymbol\pi}}\left[O\left(\varepsilon\cdot2^{2\mathbf D/3}\right)^{|\mathbf I|}\right]&=\sum_{D=0}^\infty\Pr_{\vec{\boldsymbol\pi}}[\mathbf D=D]\cdot O\left(\varepsilon\cdot2^{2D/3}\right)^{|\mathbf I|}\\
&\le2^t\cdot\sum_{D=0}^\infty O\left(\varepsilon\cdot2^{-D/3}\right)^{|\mathbf I|}\\
&\le2^t\cdot O\left(\sum_{D=0}^\infty\varepsilon\cdot2^{-D/3}\right)^{|\mathbf I|}\\
&=2^t\cdot O(\varepsilon)^{|\mathbf I|}.
\end{align*}
Therefore, our overall bound is given by
\begin{align*}
\Corr(g,h^{\oplus t})&\le\mathbb E_{\mathbf I}[2^t\cdot O(\varepsilon)^{|\mathbf I|}]=2^t\cdot\sum_{I\subseteq[t]}\Pr[\mathbf I=I]\cdot O(\varepsilon)^{|I|}\\
&=2^t\cdot\sum_{I\subseteq[t]}(1-\varepsilon)^{|I|}\cdot\varepsilon^{t-|I|}\cdot O(\varepsilon)^{|I|}\\
&\le O(\varepsilon)^t.
\end{align*}
\end{proof}

\begin{thebibliography}{99}
\bibitem[1]{source1} M. Ajtai. $\Sigma^1_1$ -formulae on finite structures. Ann. Pure Appl. Logic, 24(1):1–48, 1983. doi: 10.1016/0168-0072(83)90038-6.
\bibitem[2]{source2} Eric Allender. A note on the power of threshold circuits. In Proceedings of the 30th Symposium on Foundations of Computer Science (FOCS), pages 580–584, 1989. doi:10.1109/SFCS.1989. 63538.
\bibitem[3]{source3} Eric Allender and Vivek Gore. A uniform circuit lower bound for the permanent. SIAM Journal on Computing, 23(5):1026–1049, 1994. doi:10.1137/S0097539792233907.
\bibitem[4]{source4} Eric Allender and Ulrich Hertrampf. Depth reduction for circuits of unbounded fan-in. Inform. and Comput., 112(2):217–238, 1994. doi:10.1006/inco.1994.1057.
\bibitem[5]{source5} Andris Ambainis, Robert Špalek, and Ronald de Wolf. A new quantum lower bound method, with applications to direct product theorems and time-space tradeoffs. Algorithmica, 55(3):422– 461, 2009. doi:10.1007/s00453-007-9022-9.
\bibitem[6]{source6} James Aspnes, Richard Beigel, Merrick Furst, and Steven Rudich. The expressive power of voting polynomials. Combinatorica, 14(2):135–148, 1994. doi:10.1007/BF01215346.
\bibitem[7]{source7} László Babai. Random oracles separate PSPACE from the polynomial-time hierarchy. Inform. Process. Lett., 26(1):51–53, 1987. doi:10.1016/0020-0190(87)90036-6.
\bibitem[8]{source8} Louay M. J. Bazzi. Polylogarithmic independence can fool DNF formulas. SIAM J. Comput., 38(6):2220–2272, 2009. doi:10.1137/070691954.
\bibitem[9]{source9} Paul Beame, Russell Impagliazzo, and Srikanth Srinivasan. Approximating AC0 by small height decision trees and a deterministic algorithm for \#AC0 SAT . In 27th Conference on Computational Complexity (CCC), pages 117–125, 2012. doi:10.1109/CCC.2012.40.
\bibitem[10]{source10} Richard Beigel, Nick Reingold, and Daniel A. Spielman. The perceptron strikes back. In Proceedings of the 6th Annual Structure in Complexity Theory Conference (SCT), pages 286–291, 1991. doi:10.1109/SCT.1991.160270.
\bibitem[11]{source11} Richard Beigel and Jun Tarui. On ACC. Comput. Complexity, 4(4):350–366, 1994. doi: 10.1007/BF01263423.
\bibitem[12]{source12} Yosi Ben-Asher and Ilan Newman. Decision trees with and, or queries. In Proceedings of the 10th Conference on Structure in Complexity Theory (SCT), pages 74–81, 1995. doi: 10.1109/SCT.1995.514729.
\bibitem[13]{source13} Shalev Ben-David and Robin Kothari. Randomized query complexity of sabotaged and composed functions. Theory Comput., 14:Paper No. 5, 27, 2018. doi:10.4086/toc.2018. v014a005.
\bibitem[14]{source14} Eric Blais and Joshua Brody. Optimal Separation and Strong Direct Sum for Randomized Query Complexity. In Proceedings of the 34th Computational Complexity Conference (CCC), pages 29:1–29:17, 2019. doi:10.4230/LIPIcs.CCC.2019.29.
\bibitem[15]{source15} Ravi B. Boppana. The average sensitivity of bounded-depth circuits. Information Processing Letters, 63(5):257–261, 1997. doi:10.1016/S0020-0190(97)00131-2.
\bibitem[16]{source16} Mark Braverman. Polylogarithmic independence fools AC 0 circuits. Journal of the ACM, 57(5), 2010.
\bibitem[17]{source17} Joshua Brody, Jae Tak Kim, Peem Lerdputtipongporn, and Hariharan Srinivasulu. A strong XOR lemma for randomized query complexity, 2020. arXiv:2007.05580.
\bibitem[18]{source18} Eshan Chattopadhyay, Pooya Hatami, Kaave Hosseini, Shachar Lovett, and David Zuckerman. XOR lemmas for resilient functions against polynomials. In Proceedings of the 52nd Symposium on Theory of Computing (STOC), pages 234–246, 2020. doi:10.1145/3357713.3384242.
\bibitem[19]{source19} Lijie Chen. New Lower Bounds and Derandomization for ACC, and a Derandomization- Centric View on the Algorithmic Method. In 14th Innovations in Theoretical Computer Science Conference (ITCS), pages 34:1–34:15, 2023. doi:10.4230/LIPIcs.ITCS.2023.34.
\bibitem[20]{source20} Lijie Chen, Zhenjian Lu, Xin Lyu, and Igor C. Oliveira. Majority vs. approximate linear sum and average-case complexity below NC1 . In 48th International Colloquium on Automata, Languages, and Programming (ICALP), pages 51:1–51:20, 2021. doi:10.4230/LIPIcs.ICALP.2021.51.       
\bibitem[21]{source21} Lijie Chen and Xin Lyu. Inverse-exponential correlation bounds and extremely rigid matrices from a new derandomized XOR lemma. In Proceedings of the 53rd Symposium on Theory of Computing (STOC), pages 761–771, 2021. doi:10.1145/3406325.3451132.
\bibitem[22]{source22} Lijie Chen, Xin Lyu, and R. Ryan Williams. Almost-everywhere circuit lower bounds from non-trivial derandomization. In 61st Symposium on Foundations of Computer Science (FOCS), pages 1–12, 2020. doi:10.1109/FOCS46700.2020.00009.
\bibitem[23]{source23} Lijie Chen and Hanlin Ren. Strong average-case circuit lower bounds from nontrivial derandomization. SIAM Journal on Computing, 51(3):STOC20–115–STOC20–173, 2022. doi:10.1137/20M1364886.
\bibitem[24]{source24} Shiteng Chen and Periklis A. Papakonstantinou. Depth reduction for composites. SIAM J. Comput., 48(2):668–686, 2019. doi:10.1137/17M1129672.
\bibitem[25]{source25} Yeyuan Chen, Yizhi Huang, Jiatu Li, and Hanlin Ren. Range avoidance, remote point, and hard partial truth table via satisfying-pairs algorithms. In Proceedings of the 55th Symposium on Theory of Computing (STOC), pages 1058–1066, 2023. doi:10.1145/3564246.3585147.
\bibitem[26]{source26} Andrew Drucker. Improved direct product theorems for randomized query complexity. Comput. Complexity, 21(2):197–244, 2012. doi:10.1007/s00037-012-0043-7.
\bibitem[27]{source27} Yuval Filmus. Smolensky’s lower bound. Unpublished, 2010. URL: https://yuvalfilmus.cs. technion.ac.il/Manuscripts/Smolensky.pdf.
\bibitem[28]{source28} Merrick Furst, James B. Saxe, and Michael Sipser. Parity, circuits, and the polynomial-time hierarchy. Math. Systems Theory, 17(1):13–27, 1984. doi:10.1007/BF01744431.
\bibitem[29]{source29} Oded Goldreich, Noam Nisan, and Avi Wigderson. On Yao’s XOR-lemma. In Studies in complexity and cryptography, volume 6650 of Lecture Notes in Comput. Sci., pages 273–301. Springer, Heidelberg, 2011. doi:10.1007/978-3-642-22670-0\_23.
\bibitem[30]{source30} Shafi Goldwasser, Dan Gutfreund, Alexander Healy, Tali Kaufman, and Guy N. Rothblum. Verifying and decoding in constant depth. In Proceedings of the 39th Symposium on Theory of Computing (STOC), pages 440–449, 2007. doi:10.1145/1250790.1250855.
\bibitem[31]{source31} Aryeh Grinberg, Ronen Shaltiel, and Emanuele Viola. Indistinguishability by adaptive procedures with advice, and lower bounds on hardness amplification proofs. In Proceedings of the 59th Symposium on Foundations of Computer Science (FOCS), pages 956–966, 2018. doi:10.1109/FOCS.2018.00094.
\bibitem[32]{source32} Dan Gutfreund and Guy N. Rothblum. The complexity of local list decoding. In Proceedings of the 12th International Conference on Randomization and Computation (RANDOM), pages 455–468, 2008. doi:10.1007/978-3-540-85363-3\_36.
\bibitem[33]{source33} András Hajnal, Wolfgang Maass, Pavel Pudlák, Mario Szegedy, and György Turán. Threshold circuits of bounded depth. Journal of Computer and System Sciences, 46(2):129–154, 1993. doi:10.1016/0022-0000(93)90001-D.
\bibitem[34]{source34} Prahladh Harsha and Srikanth Srinivasan. On polynomial approximations to AC0 . Random Structures Algorithms, 54(2):289–303, 2019. doi:10.1002/rsa.20786.
\bibitem[35]{source35} Johan Håstad. Almost optimal lower bounds for small depth circuits. In Proceedings of the 18th Symposium on Theory of Computing (STOC), pages 6–20, 1986. doi:10.1145/12130.12132.
\bibitem[36]{source36} Johan Håstad. Computational limitations for small depth circuits. PhD thesis, Massachusetts Institute of Technology, 1986.
\bibitem[37]{source37} Johan Håstad. A slight sharpening of LMN. Journal of Computer and System Sciences, 63(3):498–508, 2001. doi:10.1006/jcss.2001.1803.
\bibitem[38]{source38} Johan Håstad. On the correlation of parity and small-depth circuits. SIAM J. Comput., 43(5):1699–1708, 2014. doi:10.1137/120897432.
\bibitem[39]{source39} Johan Håstad, Benjamin Rossman, Rocco A. Servedio, and Li-Yang Tan. An average-case depth hierarchy theorem for Boolean circuits. J. ACM, 64(5):Art. 35, 27, 2017. doi:10.1145/3095799.
\bibitem[40]{source40} Pooya Hatami, William M. Hoza, Avishay Tal, and Roei Tell. Depth-d threshold circuits vs. depth-(d + 1) and-or trees. In Proceedings of the 55th Symposium on Theory of Computing (STOC), pages 895–904, 2023. Full version: https://eccc.weizmann.ac.il/report/2022/ 087/. doi:10.1145/3564246.3585216.   
\bibitem[41]{source41} William M. Hoza. A technique for hardness amplification against AC0 . https://eccc. weizmann.ac.il/report/2023/176/, 2023.
\bibitem[42]{source42} Xuangui Huang and Emanuele Viola. Average-case rigidity lower bounds. In Proceedings of the 16th International Computer Science Symposium in Russia (CSR), pages 186–205, 2021. doi:10.1007/978-3-030-79416-3\_11.
\bibitem[43]{source43} Russell Impagliazzo. Hard-core distributions for somewhat hard problems. In 36th Symposium on Foundations of Computer Science (FOCS), pages 538–545, 1995. doi:10.1109/SFCS.1995. 492584.
\bibitem[44]{source44} Russell Impagliazzo, William Matthews, and Ramamohan Paturi. A satisfiability algorithm for AC0 . In Proceedings of the 23rd Symposium on Discrete Algorithms (SODA), pages 961–972, 2012. doi:10.1137/1.9781611973099.77.
\bibitem[45]{source45} Russell Impagliazzo, Ran Raz, and Avi Wigderson. A direct product theorem. In Proceedings of 9th Annual Conference on Structure in Complexity Theory (SCT), pages 88–96, 1994. doi:10.1109/SCT.1994.315814.
\bibitem[46]{source46} Rahul Jain, Hartmut Klauck, and Miklos Santha. Optimal direct sum results for deterministic and randomized decision tree complexity. Inform. Process. Lett., 110(20):893–897, 2010. doi:10.1016/j.ipl.2010.07.020.
\bibitem[47]{source47} Jeff Kahn, Gil Kalai, and Nathan Linial. The influence of variables on boolean functions. In Proceedings of the 29th Symposium on Foundations of Computer Science (FOCS), pages 68–80, 1988. doi:10.1109/SFCS.1988.21923.
\bibitem[48]{source48} Hartmut Klauck, Robert Špalek, and Ronald de Wolf. Quantum and classical strong direct product theorems and optimal time-space tradeoffs. SIAM J. Comput., 36(5):1472–1493, 2007. doi:10.1137/05063235X.
\bibitem[49]{source49} Adam R. Klivans. On the derandomization of constant depth circuits. In Proceedings of the 5th International Workshop on Randomization and Approximation Techniques in Computer Science (RANDOM), pages 249–260, 2001. doi:10.1007/3-540-44666-4\_28.
\bibitem[50]{source50} Swastik Kopparty. Lecture 4: AC0 lower bounds and pseudorandomness. Scribe notes by Jason Perry and Brian Garnett, 2013. URL: https://sites.math.rutgers.edu/\textasciitilde{}sk1233/courses/ topics-S13/lec4.pdf.
\bibitem[51]{source51} Swastik Kopparty and Srikanth Srinivasan. Certifying polynomials for AC0 [$\oplus$] circuits, with applications to lower bounds and circuit compression. Theory of Computing, 14(12):1–24, 2018. doi:10.4086/toc.2018.v014a012.
\bibitem[52]{source52} Troy Lee and Jérémie Roland. A strong direct product theorem for quantum query complexity. Comput. Complexity, 22(2):429–462, 2013. doi:10.1007/s00037-013-0066-8.
\bibitem[53]{source53} L. A. Levin. One way functions and pseudorandom generators. Combinatorica, 7(4):357–363, 1987. doi:10.1007/BF02579323.
\bibitem[54]{source54} Fu Li and David Zuckerman. Improved extractors for recognizable and algebraic sources. In Proceedings of the 23rd International Conference on Randomization and Computation (RANDOM), pages 72:1–72:22, 2019. doi:10.4230/LIPIcs.APPROX-RANDOM.2019.72.
\bibitem[55]{source55} Nathan Linial, Yishay Mansour, and Noam Nisan. Constant depth circuits, Fourier transform, and learnability. Journal of the ACM, 40(3):607–620, 1993. doi:10.1145/174130.174138.
\bibitem[56]{source56} Noam Nisan, Steven Rudich, and Michael Saks. Products and help bits in decision trees. SIAM J. Comput., 28(3):1035–1050, 1999. doi:10.1137/S0097539795282444.
\bibitem[57]{source57} Noam Nisan and Avi Wigderson. Hardness vs. randomness. J. Comput. System Sci., 49(2):149– 167, 1994. doi:10.1016/S0022-0000(05)80043-1.
\bibitem[58]{source58} Ryan O’Donnell and Karl Wimmer. Approximation by DNF: examples and counterexamples. In Proceedings of the 34th International Colloquium on Automata, Languages and Programming (ICALP), pages 195–206, 2007. doi:10.1007/978-3-540-73420-8\_19.
\bibitem[59]{source59} Alexander Razborov. A simple proof of Bazzi’s theorem. ACM Transactions on Computation Theory, 1(1), 2009. doi:10.1145/1490270.1490273.       
\bibitem[60]{source60} Alexander A. Razborov. Lower bounds on the size of constant-depth networks over a complete basis with logical addition. Mathematical Notes of the Academy of Science of the USSR, 41(4):333–338, 1987. doi:10.1007/BF01137685.
\bibitem[61]{source61} Ronen Shaltiel. Towards proving strong direct product theorems. Comput. Complexity, 12(1-2):1–22, 2003. doi:10.1007/s00037-003-0175-x.
\bibitem[62]{source62} Ronen Shaltiel. Is it possible to improve Yao’s XOR lemma using reductions that exploit the efficiency of their oracle? Comput. Complexity, 32(1):Paper No. 5, 47, 2023. doi: 10.1007/s00037-023-00238-9.
\bibitem[63]{source63} Ronen Shaltiel and Emanuele Viola. Hardness amplification proofs require majority. SIAM J. Comput., 39(7):3122–3154, 2010. doi:10.1137/080735096.
\bibitem[64]{source64} Alexander A. Sherstov. Strong direct product theorems for quantum communication and query complexity. SIAM J. Comput., 41(5):1122–1165, 2012. doi:10.1137/110842661.
\bibitem[65]{source65} Michael Sipser. Borel sets and circuit complexity. In Proceedings of the 15th Symposium on Theory of Computing, pages 61–69, 1983. doi:10.1145/800061.808733.
\bibitem[66]{source66} Roman Smolensky. Algebraic methods in the theory of lower bounds for Boolean circuit complexity. In Proceedings of the 19th Symposium on Theory of Computing (STOC), pages 77–82, 1987. doi:10.1145/28395.28404.
\bibitem[67]{source67} Roman Smolensky. On representations by low-degree polynomials. In Proceedings of 34th Annual Symposium on Foundations of Computer Science (FOCS), pages 130–138, 1993. doi:10.1109/SFCS.1993.366874.
\bibitem[68]{source68} Robert Špalek. The multiplicative quantum adversary. In Proceedings of the 23rd Conference on Computational Complexity (CCC), pages 237–248, 2008. doi:10.1109/CCC.2008.9.
\bibitem[69]{source69} Avishay Tal. Tight bounds on the fourier spectrum of AC0. In Proceedings of the 32nd Computational Complexity Conference (CCC), pages 15:1–15:31, 2017. doi:10.4230/LIPIcs. CCC.2017.15.
\bibitem[70]{source70} Jun Tarui. Probabilistic polynomials, AC 0 functions and the polynomial-time hierarchy. Theoretical Computer Science, 113(1):167–183, 1993. doi:10.1016/0304-3975(93)90214-E.
\bibitem[71]{source71} Roei Tell. On implications of better sub-exponential lower bounds for AC 0 . https://sites. google.com/site/roeitell/Expositions, 2020.
\bibitem[72]{source72} Seinosuke Toda. PP is as hard as the polynomial-time hierarchy. SIAM J. Comput., 20(5):865– 877, 1991. doi:10.1137/0220053.
\bibitem[73]{source73} Leslie G. Valiant. Graph-theoretic arguments in low-level complexity. In Proceedings of the 6th Symposium on Mathematical Foundations of Computer Science (MFCS), pages 162–176, 1977. doi:10.1007/3-540-08353-7\_135.
\bibitem[74]{source74} Emanuele Viola. The complexity of hardness amplification and derandomization. PhD thesis, Harvard University, 2006.
\bibitem[75]{source75} Emanuele Viola. On the power of small-depth computation. Found. Trends Theor. Comput. Sci., 5(1):1–72, 2009. doi:10.1561/0400000033.
\bibitem[76]{source76} Emanuele Viola. Selected challenges in computational lower bounds. SIGACT News, 48(1):39– 45, March 2017. doi:10.1145/3061640.3061648.
\bibitem[77]{source77} Emanuele Viola. New lower bounds for probabilistic degree and ac0 with parity gates. https://eccc.weizmann.ac.il/report/2020/015/, 2020.
\bibitem[78]{source78} Ryan Williams. Nonuniform acc circuit lower bounds. J. ACM, 61(1), January 2014. doi: 10.1145/2559903.
\bibitem[79]{source79} Andrew C. Yao. Theory and application of trapdoor functions. In Proceedings of the 23rd Symposium on Foundations of Computer Science (FOCS), pages 80–91, 1982. doi:10.1109/ SFCS.1982.45.
\bibitem[80]{source80} Andrew Chi-Chih Yao. Separating the polynomial-time hierarchy by oracles. In 26th Symposium on Foundations of Computer Science (FOCS), pages 1–10, 1985. doi:10.1109/SFCS.1985.49.
\bibitem[81]{source81} Andrew Chi-Chih Yao. On ACC and threshold circuits. In Proceedings of the 31st Symposium on Foundations of Computer Science (FOCS), pages 619–627, 1990. doi:10.1109/FSCS.1990. 89583.
\end{thebibliography}
\end{document}
